\chapter{Callouts}

Callouts follow Obsidian's shape: a tinted body, a saturated left bar, and an
icon plus title in the callout's own colour. They break across pages, so a
long one will not wreck the layout.

\section{The stock set}

\begin{note}
The title defaults to the callout's name. Pass an argument to override it:
\texttt{\textbackslash begin\{note\}[Something else]}.
\end{note}

\begin{tip}[Use the right one]
Fourteen types ship by default. Resist using more than three or four in a
single document — the whole point is that a callout stands out, and that
stops being true once every paragraph is one.
\end{tip}

\begin{warning}
Callouts nest badly inside tables and floats. If you need one there, put the
content in the callout and the callout in the flow, not the other way round.
\end{warning}

\begin{caution}[Data loss]
\texttt{git reset --hard} discards uncommitted work with no undo. Check
\texttt{git stash list} before you reach for it.
\end{caution}

\begin{example}[Adding your own]
One line in \texttt{styles/callouts.sty}:

\vspace{4pt}
\texttt{\textbackslash newcallout\{aside\}\{grape\}\{\textbackslash faIcon\{feather\}\}\{Aside\}}
\end{example}

\begin{question}
Should this be a callout, a footnote, or a margin note? If it interrupts the
argument, callout. If it supports it, footnote. If it comments on it, margin.
\end{question}

\begin{success}
All fourteen render, break, and nest one level deep.
\end{success}

\begin{quotebox}[Bringhurst]
The typographic page is a map of the mind. It is also a map of the society
that produced it.
\end{quotebox}

\section{Untitled asides}

When you want the tint and the bar but no header line, use \texttt{aside}:

\begin{aside}[teal]
No icon, no title — just a marked passage. Good for translations, source
notes, or a definition you want set apart without announcing itself.
\end{aside}

\section{Everything else}

\begin{abstract}
This chapter demonstrates callouts. That is all it does.
\end{abstract}

\begin{info}
Colours come from \texttt{vars/colors.tex}. Change one there and every callout
using it follows.
\end{info}

\begin{bug}[Known]
\texttt{breakable} and \texttt{parbox=false} disagree about the last line's
depth. It is off by a fraction of a point and nobody has ever noticed.
\end{bug}

\begin{idea}
A callout whose title is a question and whose body is the answer reads better
than most FAQ formatting.
\end{idea}

\begin{todo}
\begin{tasks}
  \task[done] Fourteen callout types
  \task[done] Theorem boxes with per-chapter numbering
  \task Sidenote-style callouts in the outer margin
  \task A dark-paper variant
\end{tasks}
\end{todo}

\begin{failure}
Do not put a callout inside a callout inside a callout. It compiles. It
should not.
\end{failure}

\section{Exercises}

An \texttt{exercise} is a two-part box: everything before \verb|\tcblower| is
the question, everything after it is the solution.

\begin{exercise}[Telescoping sums]
Show that $\sum_{k=1}^{n} \frac{1}{k(k+1)} = \frac{n}{n+1}$ for every
$n \geq 1$.
\tcblower
Write $\frac{1}{k(k+1)} = \frac{1}{k} - \frac{1}{k+1}$. Summing, every term
cancels against its neighbour except the first and the last, leaving
$1 - \frac{1}{n+1} = \frac{n}{n+1}$.
\end{exercise}

The title is optional; without one the box is just numbered.

\begin{exercise}
Prove that there is no largest prime.
\tcblower
Suppose $p_1, \dots, p_n$ were all of them. Then $N = p_1 \cdots p_n + 1$ has
a prime factor, and it is none of the $p_i$, since each leaves remainder 1.
\end{exercise}

\subsection{Leaving room to answer}

When the solution belongs to the reader rather than the book, use
\verb|\workspace| — a blank box on the same dot grid as the cover.

\begin{exercise}[For the reader]
Find all integers $n$ for which $n^2 + 1$ divides $n^3 + 1$.
\end{exercise}

\workspace

It takes any \texttt{tcolorbox} key, so the height and the label are
adjustable:

\workspace[height=22mm, wslabel=Working]

\begin{note}[Where the numbering comes from]
Exercises count on their own chapter-local counter, so they are numbered
independently of the theorem boxes in \cref{cha:typography}.
\end{note}
